% ============================================================
% CHAPTER 2 — The Development Environment
% ============================================================
\chapter{The Development Environment: VS~Code, Copilot, and Claude~Code}\index{VS Code}\index{GitHub Copilot}\index{Claude Code}

\section*{What You Will Learn}
By the end of this chapter you will have:
\begin{itemize}
  \item VS Code installed and configured for 0-code development
  \item GitHub Copilot activated and working (or Claude Code)
  \item Generated your first AI-assisted code snippet
  \item Understood the three main ways to interact with AI in VS~Code
\end{itemize}

% ---------------------------------------------------------
\section{Installing Visual Studio Code}

Visual Studio Code (VS~Code) is Microsoft's free code editor. It is the de facto industry standard and, more importantly for this book, the place where GitHub Copilot and Claude Code operate most effectively.

\begin{versionbox}
The instructions in this chapter were verified with VS~Code~1.96+, GitHub Copilot Chat~v0.25+, and Claude Code~v1.x. These interfaces evolve rapidly. If a button or menu changes location, rely on the underlying method rather than the exact screenshot.
\end{versionbox}

\begin{praticabox}[title={Hands-on --- Installing VS Code}]
\textbf{Windows:}
\begin{enumerate}
  \item Go to \url{https://code.visualstudio.com}
  \item Download the installer for Windows
  \item Run the installer and enable the system integration options
\end{enumerate}

\textbf{macOS:}
\begin{lstlisting}[language=bash]
brew install --cask visual-studio-code
\end{lstlisting}

\textbf{Linux (Ubuntu/Debian):}
\begin{lstlisting}[language=bash]
sudo snap install code --classic
\end{lstlisting}
\end{praticabox}

\begin{checkpointbox}
Open VS~Code. If you can see the welcome screen, your base environment is ready.
\end{checkpointbox}

% ---------------------------------------------------------
\section{Essential Extensions}

VS~Code is powerful on its own, but its real strength comes from extensions.

\begin{praticabox}[title={Hands-on --- Installing the core extensions}]
Open the Extensions panel (\texttt{Ctrl+Shift+X} on Windows/Linux, \texttt{Cmd+Shift+X} on macOS) and install the following:
\end{praticabox}

\begin{center}
\begin{tabularx}{\textwidth}{l l X}
\toprule
\textbf{Extension} & \textbf{ID} & \textbf{Purpose} \\
\midrule
GitHub Copilot & \texttt{GitHub.copilot} & Main AI engine \\
GitHub Copilot Chat & \texttt{GitHub.copilot-chat} & Chat and Agent Mode \\
ESLint & \texttt{dbaeumer.vscode-eslint} & JavaScript/TypeScript linting \\
Prettier & \texttt{esbenp.prettier-vscode} & Automatic formatting \\
Thunder Client & \texttt{rangav.vscode-thunder-client} & REST API testing \\
Flutter & \texttt{Dart-Code.flutter} & Flutter development (Part~IV) \\
Prisma & \texttt{Prisma.prisma} & Prisma schema support \\
\bottomrule
\end{tabularx}
\end{center}

\begin{tipbox}
You do not need every extension immediately. For the first part of the book, Copilot plus the JavaScript tooling is sufficient.
\end{tipbox}

% ---------------------------------------------------------
\section{Configuring GitHub Copilot}

GitHub Copilot is GitHub's integrated AI assistant inside VS~Code.

\begin{praticabox}[title={Hands-on --- Activating Copilot}]
\begin{enumerate}
  \item Create a GitHub account if you do not already have one.
  \item Choose a Copilot plan:
    \begin{itemize}
      \item \textbf{Copilot Free}: enough for the first experiments
      \item \textbf{Copilot Individual}: suitable for the whole book
      \item \textbf{Copilot Pro}: best if you want unrestricted Agent Mode
    \end{itemize}
  \item In VS~Code, open the Command Palette (\texttt{Ctrl+Shift+P}) and search for \texttt{GitHub: Sign In}
  \item Open a \texttt{.py} or \texttt{.js} file and verify that inline suggestions appear
\end{enumerate}
\end{praticabox}

\begin{costbox}
With Copilot Individual, the overall cost of following the book is typically lower than a single professional online course. Claude Code has a comparable cost profile if used with moderation.
\end{costbox}

\subsection{The three Copilot modes}

\begin{center}
\begin{tabularx}{\textwidth}{l l X X}
\toprule
\textbf{Mode} & \textbf{How to activate} & \textbf{What it does} & \textbf{When to use it} \\
\midrule
Inline completion & Automatic while typing & Suggests line-by-line code & Manual editing \\
Chat & Side panel & Explanations, snippets, file generation & Guided interaction \\
Agent Mode & Select ``Agent'' in the chat & Creates files, runs commands, iterates & \textbf{This is the 0-code mode} \\
\bottomrule
\end{tabularx}
\end{center}

\begin{warnbox}
Agent Mode is the mode you will use for most of the hands-on work in this book. It can create files, run commands, install packages, and iterate on errors with a degree of autonomy that feels like working with a junior developer.
\end{warnbox}

\begin{versionbox}
The list of available models changes frequently. As a general rule, choose the most recent and most capable model available in the Claude or GPT family.
\end{versionbox}

% ---------------------------------------------------------
\section{Alternative: Claude Code from the Terminal}

If you prefer a terminal-first workflow, you can use \textbf{Claude Code} instead.

\begin{praticabox}[title={Hands-on --- Installing Claude Code}]
\begin{lstlisting}[language=bash]
# Requires Node.js 18+
npm install -g @anthropic-ai/claude-code

# First launch and authentication
claude

# Navigate to the project folder and start
cd mio-progetto
claude
\end{lstlisting}
\end{praticabox}

\begin{center}
\begin{tabularx}{\textwidth}{l X X}
\toprule
\textbf{Aspect} & \textbf{Copilot Agent Mode} & \textbf{Claude Code} \\
\midrule
Interface & GUI inside VS~Code & Terminal \\
Learning curve & Low & Medium \\
Autonomy & High & Very high \\
Ideal for & New projects and visual workflows & Complex multi-file tasks \\
Cost & Subscription & Pay-per-use \\
\bottomrule
\end{tabularx}
\end{center}

\begin{tipbox}
All examples in this book are shown with Copilot in VS~Code, but the same context files and the same ADLC method work with Claude Code as well.
\end{tipbox}

% ---------------------------------------------------------
\section{Beyond VS Code: AI-Native IDEs in 2026}

In 2026, the tooling landscape changed rapidly. Several IDEs were designed \textbf{from the ground up} around AI orchestration.

\begin{center}
\begin{small}
\begin{tabularx}{\textwidth}{l X X X}
\toprule
\textbf{Feature} & \textbf{VS Code + Copilot} & \textbf{Cursor} & \textbf{Windsurf} \\
\midrule
Architecture & Extension-based & AI-native VS~Code fork & Standalone RAG IDE \\
Base cost & \$10/month & \$20/month & \$20/month \\
Key strength & Ecosystem and extensions & Coordinated multi-file editing & Deep repository retrieval \\
Best for & Teaching and general projects & Refactoring and complex architectures & Rapid prototyping at scale \\
\bottomrule
\end{tabularx}
\end{small}
\end{center}

\begin{notebox}
The tools will change. The \textbf{method} you learn in this book --- Context Engineering, ADLC, and disciplined verification --- remains transferable.
\end{notebox}

% ---------------------------------------------------------
\section{Your First AI-Generated Code}

\begin{praticabox}[title={Hands-on --- Generating a Python function}]
\begin{enumerate}
  \item Open an empty working folder in VS~Code
  \item Open Copilot Chat in Agent Mode
  \item Type this prompt:
\end{enumerate}
\begin{lstlisting}[language={}]
Crea un file Python chiamato calcolatrice.py che implementi una calcolatrice
 da riga di comando. Deve supportare le 4 operazioni base (+, -, *, /).
 L'utente inserisce due numeri e l'operazione da terminale.
 Gestisci la divisione per zero con un messaggio di errore chiaro.
 Aggiungi un file di test con pytest.
\end{lstlisting}
\begin{enumerate}
  \setcounter{enumi}{3}
  \item Observe what the AI creates
  \item Run the result with \texttt{python calcolatrice.py}
\end{enumerate}
\end{praticabox}

\begin{checkpointbox}
If the calculator works and the tests pass, your environment is configured correctly.
\end{checkpointbox}

% ---------------------------------------------------------
\section{How to Communicate with AI: The Basics}

\subsection*{The five rules of an effective prompt}
\begin{enumerate}
  \item \textbf{Be specific, not generic}
  \item \textbf{Show the file structure}
  \item \textbf{Define the constraints}
  \item \textbf{Specify the expected response format}
  \item \textbf{Provide project context}
\end{enumerate}

\begin{notebox}
These five rules are the simplified form of \textbf{Context Engineering}. In the next chapter we will formalize them into a reusable working method.
\end{notebox}

% ---------------------------------------------------------
\section{The Structure of a 0-Code Project}

Every project in this book follows the same center of gravity:

\begin{lstlisting}[language={}]
mio-progetto/
├── _CONTEXT.md          ← The contract with the AI
├── .copilot/
│   ├── instructions.md  ← Project-specific rules
│   └── skills/          ← Optional specialist skills
├── PROGRESS.md          ← Persistent memory across sessions
├── src/                 ← Source code generated by the AI
├── tests/               ← Tests generated by the AI
└── package.json / requirements.txt / pubspec.yaml
\end{lstlisting}

\begin{quote}
This is the center of gravity of 0-code development: \textbf{you write the context, the AI writes the code}.
\end{quote}

% ---------------------------------------------------------
\section{Recommended VS Code Settings}

\begin{praticabox}[title={Hands-on --- Recommended settings}]
Open \texttt{Open Settings (JSON)} and add:
\begin{lstlisting}[language=json]
{
  "editor.formatOnSave": true,
  "editor.defaultFormatter": "esbenp.prettier-vscode",
  "editor.minimap.enabled": false,
  "files.autoSave": "onFocusChange",
  "terminal.integrated.defaultProfile.windows": "PowerShell",
  "github.copilot.chat.agent.enabled": true
}
\end{lstlisting}
\end{praticabox}

% ---------------------------------------------------------
\section*{Summary}

\begin{center}
\begin{tabular}{l c}
\toprule
\textbf{What you did} & \textbf{Status} \\
\midrule
Installed VS~Code & \checkmark \\
Installed GitHub Copilot & \checkmark \\
Chose an AI model & \checkmark \\
Generated a first Python program & \checkmark \\
Understood the rules of a good prompt & \checkmark \\
Configured the environment & \checkmark \\
\bottomrule
\end{tabular}
\end{center}
